% 使用上游排版类, 将课程作业差异集中放在此文件。
\makeatletter
% 不插入奇偶页空白, 避免课程作业超过规定页数。
\let\cleardoublepage\clearpage
\renewcommand{\mainmatter}{\clearpage\@mainmattertrue\neu@bodymattertrue\@openrightfalse\pagenumbering{arabic}\normalsize}
\makeatother
\fancypagestyle{plain}{
  \fancyhf{}
  \fancyhead[L]{\kaishu\zihao{5}[课程名称]课程作业}
  \fancyhead[R]{\kaishu\zihao{5}\leftmark}
  \fancyfoot[C]{\zihao{-4}-\ \thepage\ -}
  \renewcommand{\headrulewidth}{0.5pt}
}
\pagestyle{plain}
% 允许页底自然留白, 避免为上下对齐拉大正文段间距。
\raggedbottom
\setcounter{tocdepth}{1}
% 将课程报告的目录压缩到单页, 保留正文原有字号与行距。
\setlength{\cftbeforechapskip}{5pt}
\setlength{\cftbeforesecskip}{0pt}
\renewcommand{\cftchapfont}{\heiti\zihao{-4}}
\renewcommand{\cftchappagefont}{\zihao{-4}}
\setlist[itemize]{nosep,leftmargin=2em}

% 保持小四号正文, 用 18 磅行距为课程报告的插图留出篇幅。
\setbodyleading{18bp}

% 课程报告采用紧凑标题间距, 正文字号和行距不变。
\ctexset{section={beforeskip=8bp,afterskip=8bp}}
% 参考文献使用五号字, 保留模板的顺序编码样式。
\AtBeginEnvironment{thebibliography}{\zihao{5}\setlength{\baselineskip}{13bp}\raggedright}

% 防止章节交界处留下很短的段落续行。
\clubpenalty=10000
\widowpenalty=10000

